\section{Colocation systems and the planner experiment}
\label{sec:appx-gap}

Colocation runs best-effort work on resources that latency-critical
serving leaves unused while preserving its \acp{SLO}
~\cite{flexllm,llmstation,harli,deltaserve}. Fine-tuning is a natural
candidate because it is throughput-oriented and delay-tolerant, and
\ac{PEFT} methods such as \ac{LoRA}~\cite{lora} train only small
adapters while reusing the serving model's resident weights. Because
GPU kernels are difficult to preempt once launched~\cite{nonpreempt},
colocation systems must both identify spare capacity and expose work
units short enough to fit inside it.

Existing inference--fine-tuning systems already decompose training for
this purpose, but make admission decisions at serving boundaries using
signals such as batch composition, iteration timing, or engine
headroom (\Cref{tab:colocation_systems}). Expert-parallel bubbles have a different
granularity: they open and close inside one \ac{MoE} layer and differ
across GPUs. Using them therefore requires both \emph{bubble-sized
work} and a \emph{per-GPU} admission decision that reveals where slack
exists and for how long. The current router assignment provides this missing signal
(\S\ref{sec:design-predict}).

\begin{table*}[t]
\centering
\small
\caption{
\textbf{Existing colocation systems reason about serving-level headroom, whereas \sys{} identifies slack per GPU inside an \ac{MoE} layer.}
}
\label{tab:colocation_systems}
\begin{tabular}{lllll}
\toprule
\emph{System} & \emph{Capacity targeted} & \emph{Work unit} & \emph{Admission scope} & \emph{Signal} \\
\midrule
\flexllm~\cite{flexllm}
& serving headroom
& token slot
& serving iteration
& batch composition \\
\llmstation~\cite{llmstation}
& decode / idle headroom
& tasklet
& serving iteration
& iteration timing \\
\harli~\cite{harli}
& decode underutilization
& SM partition
& serving iteration
& predicted latency \\
\deltaserve~\cite{deltaserve}
& serving / idle headroom
& LoRA batch
& serving iteration
& predicted latency \\
\midrule
\sys{}
& \ac{EP} bubble
& fine-tuning chunk
& GPU $\times$ \ac{MoE} layer

& router top-$k$ assignment \\
\bottomrule
\end{tabular}
\end{table*}

We first isolate the planning requirement. We replay a
Qwen3-30B-A3B serving trace (\ac{TP}$=$\ac{EP}$=$2, 128k-token
prefill) against the same 486 measured fine-tuning chunks
(median 567\,$\mu$s) while varying only the information available to
the planner. A \emph{router-blind} planner assigns every GPU a
historical layer-level slack estimate; an \emph{oracle GPU-agnostic}
planner instead receives the true mean bubble of the current layer but
still assigns the same budget to every GPU; and a \emph{router-aware}
planner estimates slack separately for each GPU from the current
routing decision. We compare their \emph{bubble utilization}, the
fraction of available bubble GPU-time occupied by training, at
approximately the same 1\% overrun rate.

\begin{table}[t]
  \centering
  \small
  \caption{
    \textbf{Per-GPU planning is required to use bubbles.}
    All planners use the same fine-grained fine-tuning chunks and operate
    near a 1\% overrun rate.}

  \label{tab:gap-planner}
  \begin{tabular}{lrr}
    \toprule
    Planner & Bubble utilization & Overrun rate \\
    \midrule
    Router-blind
      & 1.28--1.31\% & 1.04\% \\
    Oracle GPU-agnostic
      & 1.22--1.30\% & 1.04\% \\
    Router-aware
      & 42.74--45.14\% & 0.83--1.04\% \\
    \bottomrule
  \end{tabular}
\end{table}

Table~\ref{tab:gap-planner} shows that router-aware planning recovers
42.7--45.1\% of bubble capacity, compared with about 1.3\% for either
GPU-agnostic planner. Even knowing the exact \emph{average} slack of
the current layer is therefore insufficient: the planner must know
which GPU owns that slack.

Shorter work units determine how much of this per-GPU capacity can be
used. With router-aware planning, reducing the median chunk duration
from 4.13\,ms (97 chunks) to 567\,$\mu$s (486 chunks) raises bubble
utilization from 37.95--38.46\% to 42.74--45.14\%. Yet the
router-blind planner remains near 1--1.5\% under both partitions.
Thus, finer execution helps only after the planner identifies the
right GPU: chunk size determines \emph{how much} work fits, while
router state determines \emph{where} it can safely run.
